% ch39-svdq-quant-producer.tex — SVDQuant(二)：量化 producer（RFC-P2）
% 源码：cache-dit(db125215) csrc/kernels/svdq/cute/fp4/quantize/ + gemm/ + ops.h
% 公式锚点一律回指 ch38（sec:38-*），本章不重推。
\chapter{SVDQuant（二）：量化 producer——融合 LoRA-down 的激活量化}\label{ch:39}

\section{本章导读}

第~\ref{ch:38}~章的式~\ref{eq:38-svdq} 说，SVDQuant 在线侧要做两次计算：
激活量化 $\mathcal Q(\widehat X)$ 与低秩 down 投影
$D=\widehat XL_2^{\top}=X(L_2S^{-1})^{\top}$。两者读的是同一份输入
$X$，朴素写法会把它变成多条独立的 kernel 链，引入额外中间张量的 DRAM
读写。本章讲 cache-dit 的\textbf{融合 producer}：在同一个公共算子内共享
激活装载，同时产出 GEMM 的激活码、scale 与 LoRA-down 结果。当前实现有
fused 与 split-K 两种调度；前者用单个 kernel 完成，后者还需要一次 BF16
finalize，不能再把整条链概括成“永远只有一个 kernel”。

\textbf{源码区间}：cache-dit 仓库 commit \texttt{db125215} 的
\texttt{csrc/}\allowbreak\texttt{kernels/}\allowbreak\texttt{svdq/}%
\allowbreak\texttt{cute/}\allowbreak\texttt{fp4/}\allowbreak\texttt{quantize/}。

\texttt{quantize\_uni.cu} 是通用实现；\texttt{quantize\_impl.cuh} 是 rank
特化的 CuTe \texttt{TiledMMA}/\allowbreak\texttt{TiledCopy} 实现。
\texttt{quantize\_r32.cu}、\texttt{quantize\_r64.cu} 与
\texttt{quantize\_r128.cu} 分别只实例化一个 rank；公共声明在
\texttt{quantize.h}。实际选择还取决于
\texttt{csrc/kernels/svdq/}\allowbreak\texttt{ops.h}，不能仅凭文件存在就推断
默认路径。GEMM 的消费约定在
\texttt{cute/fp4/gemm/}\allowbreak\texttt{gemm\_w4a4.h}，\texttt{LoraLayout}
枚举统一放在 \texttt{svdq/runtime.h}。完整代码入口见附录~\ref{app:D}。

\textbf{架构限定}：device 分支要求
\texttt{\_\_CUDA\_ARCH\_\_ >= 1200}，否则落入 \texttt{asm volatile("trap;")}。
这与 FP4 量化指令 \texttt{cvt.rn.satfinite.e2m1x2.f32} 的硬件门槛一致
（第~\ref{ch:31}~章）。producer 的 host \texttt{require} 链本身不查询设备，
能力检查不能靠这段校验代替。

\textbf{能力边界}：输入与 \texttt{smooth} 必须 BF16 且 contiguous；
\texttt{qact} 必须 INT8、\texttt{ascales} 必须 FP8\_E4M3；
$M\equiv0\pmod{256}$，$K\equiv0\pmod{64}$。通用实现接受 rank 0 或不超过
128 的 16 的倍数，\texttt{lora\_scales} 长度为 $\max(1,\mathrm{rank}/16)$；
rank 特化入口只接受与 TU 相符的 32/64/128，scale 数分别为 2/4/8。
split-K 还要求至少 $M\times R$ 个元素的 FP32 scratch。

\textbf{验证方式}：cache-dit 的 \texttt{tests/kernels/test\_svdquant\_cute.py}
（SM120 门控，非 SM120 自动 skip；本书未复跑），本章最小测试节另给出
CPU 侧的 pack 复算思路。

\section{问题与动机}

\subsection{在线侧两次计算，一份输入}

把式~\ref{eq:38-svdq} 的在线部分展开成数据流：每个 forward，
$X$ 先被逐通道缩放平滑除一次得到 $\widehat X$，随后一路做 NVFP4 量化
（第~\ref{ch:31}~章的两级 scale：每 16 元素一个 e2m1 块配一个 ue4m3
SFA），另一路乘低秩 down 因子得到 $D$（$M\times r$ 的 bf16）。
\textbf{关键实现事实}：离线打包时 $L_2$ 已经被预乘了 $S^{-1}$
（$\widetilde L_2=L_2S^{-1}$，cache-dit 的 packing 代码写作
\texttt{lora\_down /= smooth}），所以
\begin{equation}
  D=X\widetilde L_2^{\top}=XS^{-1}L_2^{\top}=\widehat XL_2^{\top},
  \label{eq:39-d}
\end{equation}
LoRA 路径可以直接消费\textbf{原始} $X$，与理论式完全一致。这正是
第~\ref{ch:38}~章“常见坑”里“S 最容易重复乘除”的具体体现：量化路径只对
临时副本除 \texttt{smooth}，LoRA 路径继续使用原始 A fragment，两边各自
正确、互为镜像。

\subsection{朴素方案与融合方案}

朴素方案按算子逐个实现：smooth kernel 读 $X$ 写 $\widehat X$；量化
kernel 读 $\widehat X$ 写码与 SFA；LoRA-down GEMM 读 $X$（或
$\widehat X$）写 $D$；如果 GEMM 要的 D 布局与 GEMM 的输入布局不一致，
再来一个 finalize/repack kernel。SVDQuant 论文 Fig.~6 的融合原则
（第~\ref{ch:38}~章性能节引用过）要求 down projection 与 activation
quantize \textbf{共享同一份输入}。当前 fused 调度用一个 kernel 产出三项
输出；split-K 也让量化与 LoRA-down 共享 A 装载，只是额外经过内部 FP32
scratch 和 finalize：
\begin{center}\small
\begin{tabular}{@{}lll@{}}
\toprule
输出 & 形状/格式 & 消费者 \\
\midrule
\texttt{qact} & $[M,K/2]$ INT8，native 256 行交织 E2M1 pack & GEMM 的 A 操作数（G2R）\\
\texttt{ascales} & 与 native pack 对齐的 ue4m3 SFA & GEMM 的 SFA（G2R）\\
\texttt{lora\_act}（即 $D$） & $[M,R]$ row-major BF16，scale 已折入 & GEMM 的 LoRA 输入（TMA）\\
\bottomrule
\end{tabular}
\end{center}
scratch 是 split-K 的内部中间态，不是公共输出。第三项输出的布局选择就是
下一节的主角：\texttt{LoraLayout}。

\subsection{两种 LoRA 输入 layout（约定层对比）}

\texttt{svdq/runtime.h} 定义 \texttt{LoraLayout\{kNative,kRowMajor\}}，
描述的是\textbf{$D$ 以什么形态到达 GEMM}；\texttt{gemm/gemm\_w4a4.h}
L32--42 进一步给出两条消费路径：

\begin{itemize}[itemsep=1pt,topsep=2pt]
  \item \texttt{kNative}：$D$ 是 native producer 输出的 \textbf{FP32
        MMA-fragment pack}；GEMM 把它从 gmem gather 进寄存器，并在读取时
        施加 \texttt{lora\_scales[i]}（每 16 rank 一组）。$D$ 不进 TMA
        stage，LoRA 项只能在 NVFP4 mainloop 后执行（第~\ref{ch:40}~章细讲）。
  \item \texttt{kRowMajor}：$D$ 是\textbf{普通 row-major BF16}
        $[M,R]$，per-rank scale 已在 producer 里折入；GEMM 用 TMA 加载它，
        复用 A/B 的 stage 存储，因此 LoRA 项可以在 mainloop 前执行。
\end{itemize}

两个 layout 算的是同一个数学对象，但浮点舍入路径不同。
\texttt{kRowMajor} 在 producer 折入 \texttt{lora\_scales} 后把结果舍入成
BF16；\texttt{kNative} 把 FP32 的 $D$ 留到 GEMM 侧，gather 时乘完
\texttt{lora\_scales} 再转 BF16。因此“数学等价”不等于两条路径逐位相同。

本章讨论的是 \texttt{kRowMajor} 输出。公共 Python 入口仍是
\texttt{svdq\_}\allowbreak\texttt{quantize\_}\allowbreak\texttt{w4a4}，传入
\texttt{lora\_layout=}\allowbreak\texttt{"row\_major"} 后，\texttt{ops.h}
再选择 \texttt{quantize\_uni} 或某个 rank 特化入口。native producer 在
\texttt{csrc/kernels/svdq/}\allowbreak\texttt{native/}，本章只使用它的输出约定。
图~\ref{fig:39-lora-layout} 把两种形态画在一起。

\section{数学原理到变量的映射}

第~\ref{ch:38}~章的公式在本 kernel 里各有一个实体，先把映射表立起来，
后文逐段讲代码时直接引用：

\begin{center}\small
\begin{tabular}{@{}llll@{}}
\toprule
数学量 & kernel 变量 & 形状/类型 & 当前代码位置 \\
\midrule
$X$ & \texttt{input} & $[M,K]$ BF16 & \texttt{quantize\_cute}/\texttt{quantize\_uni} 参数 \\
$S$ & \texttt{smooth} & $[K]$ BF16 & \texttt{smooth\_fragment} 的输入 \\
$\widehat X=X/S$ & 临时 \texttt{PackedBF16} & fragment & 只改临时副本，保留原始 A \\
$\bar q,\Delta$（e2m1 码 + SFA） & \texttt{qact}/\texttt{ascales} & INT8 / ue4m3 & \texttt{quantize\_group} \\
$\widetilde L_2=L_2S^{-1}$ & \texttt{lora\_down} & $[K,R]$ BF16 & 逻辑 B 视图 stride $(1,R)$ \\
$D$ & \texttt{lora\_act} & $[M,R]$ BF16 & fused epilogue 或 split-K finalize \\
per-16-rank 折叠系数 & \texttt{lora\_scales} & 每 R16 一个 FP32 & 特化 \texttt{float[8]}；通用 \texttt{float[64]} \\
\bottomrule
\end{tabular}
\end{center}

两处值得停一下。第一，量化步长按第~\ref{ch:31}~章的 NVFP4 语义：
\begin{equation}
  \Delta_{block}=\min\!\left(\frac{\mathrm{amax}(16\text{ 元素块})}{6},\;448\right),
  \qquad \bar q=\mathrm{round}\!\left(x\cdot\mathrm{rcp\_approx}(\Delta)\right),
  \label{eq:39-scale}
\end{equation}
除以 6 是 e2m1 满量程，clamp 到 448 是 ue4m3（无符号 e4m3）的值域上限。
当前 producer 的若干注释写“per-K32 scale”，但 host 校验要求
\texttt{ascales} 至少容纳 $MK/16$ 字节，GEMM 约定也明确每 16 个 K 元素
共享一个 FP8\_E4M3 scale。本文按实现与消费约定采用 \textbf{per-16}，不把
寄存器重组时的 K32 面板误当成 scale 粒度。

这里还有两处舍入要分清：量化侧乘浮点 $\Delta$ 的\textbf{近似倒数}
（\texttt{rcp.approx.ftz.f32}），反量化侧乘\textbf{打包进 SFA 的 e4m3
舍入之后的 $\Delta$}。$\bar q$ 的落点按 $1/\Delta_{\text{float}}$ 算，
回读时按 $\Delta_{\text{e4m3}}$ 乘，两者相差一次 e4m3 舍入（相对误差
量级 $2^{-4}$）。这是 NVFP4 语义的固有部分，不是 kernel 的宽容。
第二，\texttt{lora\_scales} 在每 16 个 rank 列上共享一个折叠系数：
$D\leftarrow\mathrm{bf16}(\mathrm{fp32}(D)\cdot
s_{\text{lora}}[\lfloor\mathrm{col}/16\rfloor])$。fused 在协作 epilogue
中执行这一步；split-K 把它推迟到 finalize，二者都只在最终写 $D$ 时转
BF16。

\section{设计与实现}

\subsection{两套实现、两种调度与实际 dispatch}\label{sec:39-structure}

当前 producer 不是一份 kernel 的局部改版，而是两套实现共享一套输出约定。
\texttt{quantize\_uni.cu} 保留手写 \texttt{ldmatrix/mma} 的通用路径，覆盖
rank 0、非特化 rank，以及 $R=32,K\le3072$；\texttt{quantize\_impl.cuh}
用 CuTe 描述 rank 32/64/128 的特化路径。三个很短的
\texttt{quantize\_r32/64/128.cu} 分别调用 \texttt{quantize\_cute<2/4/8>}，
让每个 TU 只实例化一个 rank，编译可以并行，调参也不会把另外两个 rank 的
模板实例一起拖进来。

\texttt{ops.h} L161--189 先做两级选择。第一级按 rank 与 $K$ 选实现；第二级
用是否分配 \texttt{lora\_scratch} 选调度。有效 scratch 是 split-K 的显式
选择器，不是可随意忽略的优化提示。当前默认规则如下：

\begin{center}\small
\begin{tabular}{@{}lllll@{}}
\toprule
rank & 形状条件 & 实现 & 调度 & K chunk / A copy \\
\midrule
0 或非特化 rank & 任意 & \texttt{quantize\_uni} & fused & warp 切 K64 \\
32 & $K\le3072,M\le8192$ & \texttt{quantize\_uni} & split-K & K128，同步 \\
32 & $K\le3072,M>8192$ & \texttt{quantize\_uni} & fused & warp 切 K64 \\
32 & $K>3072,M\le8192$ & \texttt{quantize\_r32} & split-K & K128，双 stage 异步 A \\
32 & $K>3072,M>8192$ & \texttt{quantize\_r32} & fused & warp 切 K64 \\
64 & $K\le3072$ & \texttt{quantize\_r64} & split-K & K256，同步 A \\
64 & $K>3072$ & \texttt{quantize\_r64} & split-K & K128，双 stage 异步 A \\
128 & 任意 & \texttt{quantize\_r128} & split-K & K256，同步 A \\
\bottomrule
\end{tabular}
\end{center}

这些阈值是当前 PRO 5000 上调好的工程选择，不是由 CuTe 或硬件推导出的
普遍定律。直接调用 rank 入口时仍能选择 rank64/128 的 fused 实例；只是当前
\texttt{ops} 对它们总是提供 scratch，所以默认进入 split-K。还要区分两种
“split”：fused 内部的 \texttt{kRankSplits} 把 rank 列分给不同 CTA，以增加
小 $M$ 下的 grid 并行度；split-K 则把归约轴 $K$ 分给不同 CTA，并用全局
scratch 合并部分积。前者只让 \texttt{blockIdx.y==0} 写 qact/ascales，后者
的每个 $(\text{row panel},K64)$ 也只有一个激活写者。

\begin{figure}[H]
\centering
\begin{tikzpicture}[font=\scriptsize,
  box/.style={draw, rounded corners=2pt, align=center, inner sep=3pt},
  arr/.style={->, thick}]
  \node[box, fill=tkFillBlue] (in) at (0,2.0) {$X\ [M,K]$ BF16};
  \node[box, fill=tkFillGreen, text width=3.1cm] (copya) at (3.5,3.1)
    {A：TiledCopy / cp.async\\SW128 smem $\to$ LDSM\_N};
  \node[box, fill=tkFillGreen, text width=3.1cm] (copyb) at (3.5,0.9)
    {B：TiledCopy / cp.async\\SW32 smem $\to$ LDSM\_T};
  \node[box, fill=tkBlue!18] (raw) at (7.0,3.1) {原始 A fragment};
  \node[box, fill=tkPurple!18] (bf) at (7.0,0.9) {B fragment};
  \draw[arr] (in.north east) -- (copya.west);
  \draw[arr] (in.south east) -- (copyb.west);
  \draw[arr] (copya) -- (raw);
  \draw[arr] (copyb) -- (bf);
  \node[box, fill=tkOrange!22, text width=2.6cm] (quant) at (10.3,4.0)
    {临时副本除 \texttt{smooth}\\\texttt{quantize\_group}};
  \node[box, fill=tkGreen!30, text width=2.6cm] (qout) at (14.1,4.0)
    {qact + ascales\\native 256 行 pack};
  \draw[arr] (raw) -- (quant);
  \draw[arr] (quant) -- (qout);
  \node[box, fill=tkPurple!25] (mma) at (10.3,1.4)
    {TiledMMA\\$32{\times}16$ / R16 chunk\\FP32 累加};
  \draw[arr] (raw.south) -- (mma.north west);
  \draw[arr] (bf) -- (mma);
  \node[box, fill=tkFillBlue, text width=2.6cm] (fused) at (8.3,-0.6)
    {fused\\warp 切 K\\shared 归约};
  \node[box, fill=tkFillBlue, text width=3.0cm] (split) at (13.0,-0.6)
    {split-K\\K CTA $\to$ FP32 scratch\\relaxed atomic + finalize};
  \draw[arr] (mma.south west) -- (fused.north);
  \draw[arr] (mma.south east) -- (split.north);
  \node[box, fill=tkBlue!30] (D) at (10.7,-2.1)
    {$D\ [M,R]$ row-major BF16\\最终一次舍入};
  \draw[arr] (fused.south) -- (D.north west);
  \draw[arr] (split.south) -- (D.north east);
  \node[anchor=north west, align=left, text width=15.8cm, font=\tiny] at (-0.8,-3.0)
    {量化支路只修改 A 的临时副本，LoRA 支路继续消费原始 A。fused 在 CTA 内
     归约；split-K 跨 K CTA 累加 FP32 scratch，再由同一 stream 上的 finalize
     折入 R16 scale 并写 BF16。};
\end{tikzpicture}
\caption{rank 特化 producer 的数据通路。TiledMMA/TiledCopy 改变的是装载与
LoRA 计算的组织方式，qact/ascales 的 native pack 和 $D$ 的 row-major BF16
约定保持不变；fused 与 split-K 的差别集中在 $D$ 的归约位置。}
\label{fig:39-producer-flow}
\end{figure}

\subsection{rank 特化路径：用 TiledMMA/TiledCopy 组织装载}

\texttt{QuantizeTraits} 把 fused 与 split-K 共用的 CuTe 几何集中在一处。
下面 39 行是理解当前实现的入口：

\begin{lstinputlisting}[linerange={62-100},firstnumber=auto,
  title={quantize\_impl.cuh L62-100（MMA、shared layout 与 Copy atom）}]{../../../../cache-dit/csrc/kernels/svdq/cute/fp4/quantize/quantize_impl.cuh}
\end{lstinputlisting}

MMA 使用 \texttt{SM80\_16x8x16\_F32BF16BF16F32\_TN}，
\texttt{Tile<\_32,\_16,\_16>} 把逻辑计算面扩到 $32\times16$。这里的
atom 布局是 $(1,1,1)$，线程切片按 \texttt{threadIdx.x\%32} 取得，因此它是
\textbf{单 warp} 的计算描述，不是四个 warp 拼出一个 tile。一个 K64 切片
拆成四个 K16 atom；split-K 每次只保留一个 R16 的 FP32 accumulator，避免
rank128 常驻完整的 $32\times128$ 寄存器累加面。

A 的逻辑 tile 是 $32\times64$，K 维连续，shared 使用
\texttt{Layout\_K\_SW128\_Atom}；B 虽在 gmem 中是 row-major $[K,R]$，逻辑上
却按 $[N=16,K=64]$ 建立 stride $(1,R)$ 的视图，shared 使用
\texttt{Layout\_MN\_SW32\_Atom}。G2S 的 Copy atom 都是 128-bit；S2R 则分别选
\texttt{LDSM\_N} 与转置形式 \texttt{LDSM\_T}。B 的转置 copy 直接生成 BF16
MMA 所需的 K-paired fragment，不再逐标量拼凑。

同步 A 装载把 TiledCopy 的分区、shared 写、LDSM 读串成一条明确顺序：

\begin{lstinputlisting}[linerange={280-306},firstnumber=auto,
  title={quantize\_impl.cuh L280-306（A 的 G2S、LDSM 与 warp 同步）}]{../../../../cache-dit/csrc/kernels/svdq/cute/fp4/quantize/quantize_impl.cuh}
\end{lstinputlisting}

第一个 \texttt{\_\_syncwarp} 发布 shared 写，第二个防止下一轮复用 stage
时覆盖尚未读完的数据。split-K 的异步变体为 A 准备双 stage：
\texttt{cp\_async\_wait<1>} 只等待较早的 A 组，让后发射的 B 组仍在飞行；
量化完成后才用 \texttt{cp\_async\_wait<0>} 与 \texttt{\_\_syncthreads}
发布全部 B tile。$K256$ 配置使用同步 A；源码的 \texttt{static\_assert}
要求双 stage 异步 A 只用于 $kMaxSlices\le2$，即最多两个 K64 切片。

通用 \texttt{quantize\_uni.cu} 仍是另一套结构：$16\times72$ 手写 staging、
显式 \texttt{ldmatrix.x4} 与手写 \texttt{mma}。它没有被这些 Traits 替换；
当前 dispatch 之所以保留它，是因为 rank32、$K\le3072$ 的实测仍由它占优。

\subsection{量化支路：\texttt{quantize\_group} 的四个阶段}

量化支路的核心仍是一个共享的 \texttt{quantize\_group}。rank 特化路径把
\texttt{smooth\_fragment} 作用于 A 的临时副本，再将四个 K16 fragment 传入；
LoRA MMA 继续持有未经 smooth 的原始 A。一次调用覆盖 M16$\times$K64，外层
$j=0,1$ 选择两个 8 行半，每行有四个 K16 组。

\textbf{第一阶段：寄存器重排与 amax 归约}。fragment 寄存器顺序是 MMA 布局，
不是普通行布局。代码先把四个 \texttt{PackedBF16} 重组为 \texttt{input[8]}，
再对每个 K16 组求绝对值最大值；同一行组跨四个 lane，用两轮
\texttt{shfl\_xor}（mask 2、1）得到共同 amax，不需要 shared 归约。

\begin{lstinputlisting}[linerange={185-223},firstnumber=auto,
  title={quantize\_impl.cuh L185-223（寄存器重排、amax 与四个 per-16 scale）}]{../../../../cache-dit/csrc/kernels/svdq/cute/fp4/quantize/quantize_impl.cuh}
\end{lstinputlisting}

这里还要校正一处注释：L173--175 把 $j$ 称为“K32 panel”，但 L247 的
行号只随 $j$ 增加 8，两个 $j$ 的 \texttt{dst[0/4/8/12]} 写地址相同；
$j$ 实际选择行半，不是新的 K32 面板。

\textbf{第二阶段：步长与 SFA}。按式~\ref{eq:39-scale} 求 $\Delta$ 与近似
倒数 \texttt{rscale}；四个 K16 组的 $\Delta$ 经两条
\texttt{cvt.rn.satfinite.}\allowbreak\texttt{e4m3x2.f32} 打成一个
4 字节 scale 字。这里再次看到
per-16 粒度：一个 scale 字正好对应当前 K64 的四个组。

\textbf{第三阶段：e2m1 量化与跨 lane 归并}。每个 lane 把本地 BF16 对乘上
\texttt{rscale}，用 \texttt{cvt.rn.satfinite.}\allowbreak\texttt{e2m1x2.f32}
转成字节，再左移
$(\mathrm{lane}\bmod4)\times8$ 位落入字节槽。两轮 \texttt{shfl\_xor}
（mask 1、2）把四份字节贡献或在一起，得到完整的 u32；一个 u32 装 8 个
e2m1，所以它是 \textbf{K8 码字}。归并后四 lane 都有副本，实际写回只让
$\mathrm{lane}\bmod4<2$ 的两个 lane 落笔。

\begin{lstinputlisting}[linerange={224-263},firstnumber=auto,
  title={quantize\_impl.cuh L224-263（E2M1 归并与 native pack 的 qact 写出）}]{../../../../cache-dit/csrc/kernels/svdq/cute/fp4/quantize/quantize_impl.cuh}
\end{lstinputlisting}

\textbf{第四阶段：写 native 256 行交织 pack}。\texttt{qact} 不是普通行主序
$[M,K/2]$，而是把 256 行作为块，按 SM120 MMA 的 A-operand fragment 位置，
用 \texttt{base}、\texttt{local} 与 writer lane 的 \texttt{i*2} 三层偏移
排布。每个 lane 在归并后持有 8 个 u32 的副本，writer lane 0/1 分别写其中
4 个，\texttt{dst[0/4/8/12]} 对应四个 K16 组。

SFA 由 $\mathrm{lane}\bmod4=0$ 的 lane 写出：\texttt{base*32+scale\_lane}
选中与 qact 对应的位置，一个 u32 装四个 per-16 FP8 scale。通用与 rank 特化
producer 都遵守这份布局；qact 与 native 逐字节一致，ascales 也按同一映射
写出，但现有测试仅比较排序后的值集合，不能把它当作顺序逐位断言。第~\ref{ch:40}~章的 GEMM 因而不需要激活 repack，
直接用 G2R 将 A/SFA 装进寄存器，不经过 TMA 或 shared staging。

\begin{figure}[H]
\centering
\begin{tikzpicture}[font=\scriptsize,
  cell/.style={draw, minimum height=0.44cm, inner sep=1.5pt}]
  % (a) 256 行块面板
  \node[anchor=west, font=\small\bfseries] at (0,4.3)
    {(a) 一个 256 行块的 \texttt{qact} 面板};
  \foreach \c/\lab in {0/{k 块 0}, 1/{k 块 1}, 2/{$\cdots$}}
    \node[font=\tiny] at ({1.4+1.78*\c},3.92) {\lab};
  \foreach \r in {0,1,2,3,4,5,6,7}{
    \node[font=\tiny, anchor=east] at (-0.08,{3.45-0.48*\r}) {panel \r};
    \foreach \c in {0,1,2}
      \node[cell, fill=tkGreen!22, minimum width=1.58cm]
        at ({1.4+1.78*\c},{3.45-0.48*\r}) {32 行};
  }
  % (b) lane 归并与写回
  \node[anchor=west, font=\small\bfseries] at (8.5,4.3)
    {(b) 一个 M16$\times$K64 量化组的写回};
  \foreach \l in {0,1,2,3}{
    \node[font=\tiny, anchor=east] at (8.25,{3.45-0.62*\l}) {lane \l};
    \node[cell, fill=tkOrange!22, minimum width=5.4cm]
      at (12.1,{3.45-0.62*\l}) {8 个 K8 码字的完整副本};
  }
  \node[font=\tiny, anchor=west, text=tkOrange!75!black, text width=6.6cm] at (8.3,0.42)
    {仅 lane 0/1 各写 4 个 u32；lane 2/3 保留归并副本但不落笔};
  % 注记
  \node[anchor=north west, align=left, text width=15.8cm, font=\tiny] at (0,-0.25)
    {(a) 每一列是一个 K64 块（每行占其中 32 B 码），每格是 32 行
     $\times$ 32 B $=1$ KiB；行是 8 个 32 行面板，面板号
     $p=\lfloor(\mathrm{row}\bmod256)/32\rfloor$ 进入 \texttt{base}，不在
     \texttt{local}。面板内再由 \texttt{local}$=16(\mathrm{row}\bmod 8)
     +\lfloor(\mathrm{row}\bmod 16)/8\rfloor+128\lfloor(\mathrm{row}\bmod 32)/16\rfloor$
     交织出 32 个槽位；\texttt{base} 同时按 $\lfloor\mathrm{row}/256\rfloor$、
     k 块号与 $p$ 推进（\texttt{quantize\_impl.cuh} L249--252）。顺序与 SM120
     MMA 的 A-operand fragment 位置一致，所以第~\ref{ch:40}~章的 G2R 能逐字
     直读。(b) 的两轮 \texttt{shfl\_xor}（mask 1、2）把四片字节或成 u32
     （K8 码字）；每行 8 个 u32 拼满 K64，即 32 B 码。};
\end{tikzpicture}
\caption{native 激活 pack 的两级结构。(a) 行维：256 行一块、8 个 32 行面板；
(b) 写回：四个 lane 先获得相同的 K8 码字副本，再由其中两个 lane 落笔。}
\label{fig:39-native-pack}
\end{figure}

\subsection{LoRA 支路：从标量 B 加载到 \texttt{LDSM\_T}}\label{sec:39-lora-b}

通用实现最能说明 B 操作数的布局陷阱。它用 \texttt{load\_lora\_frags}
把原始 $X$ 与 \texttt{lora\_down} 累乘进 FP32 \texttt{acc}，其中
\texttt{load\_lora\_pair} 必须显式跨行取值：

\begin{lstinputlisting}[linerange={87-113},firstnumber=auto,
  title={quantize\_uni.cu L87-113（通用路径的 BLayout 与标量加载）}]{../../../../cache-dit/csrc/kernels/svdq/cute/fp4/quantize/quantize_uni.cu}
\end{lstinputlisting}

\texttt{SM80\_16x8x16\_F32BF16BF16F32\_TN} 的 BLayout 是
$((4,8),(2,2)):((16,1),(8,64))$：每个 32-bit 寄存器装的是同一 N/rank
列上、K 方向相邻的两个 BF16。而 \texttt{lora\_down} 是 row-major
$[K,R]$，这两个值在内存里相隔 $R$ 个元素。若直接读
\texttt{ptr[0],ptr[1]}，拿到的是相邻 rank 列的同 K 值，会产生静默的
even/odd 列置换。通用路径的两条标量 load 既对上布局，也避开奇数 rank
列的 4 字节对齐假设。uni 的 rank128 是例外：L459--468 改用四次对齐
32-bit load 加 \texttt{movmatrix.trans}，得到的 K-paired 寄存器顺序与标量
版本相同；rank 不超过 64 时仍保留标量路径。

rank 特化路径不再复制这套标量过程。\texttt{issue\_b\_warp} 或
\texttt{issue\_b\_cta} 先用 TiledCopy 把逻辑 $[N,K]$ 视图放进 SW32 shared
布局；随后 \texttt{mma\_b\_chunk} 通过 \texttt{LDSM\_T} 形成四个 K16 的
B fragment，并逐个调用 \texttt{cute::gemm}：

\begin{lstinputlisting}[linerange={405-423},firstnumber=auto,
  title={quantize\_impl.cuh L405-423（LDSM\_T 形成 B fragment 并交给 TiledMMA）}]{../../../../cache-dit/csrc/kernels/svdq/cute/fp4/quantize/quantize_impl.cuh}
\end{lstinputlisting}

这里的改进不只是少了几次 load 指令。布局转换由 shared layout、Copy atom
与 MMA partition 共同表达，B 的坐标映射不再需要散落在手写索引公式里。
fused 在 warp 内双缓冲 B stage；split-K 则由 128 线程 copy team 预取 CTA
所需的全部 R16$\times$K64 tile，量化完成后再统一进入 LoRA MMA。

\subsection{fused：CTA 内归约与一次性舍入}

fused 调度中，一个 CTA 拥有 32 行，warp 沿 K64 切片循环。K 循环结束后，
每个 warp 持有自己那段 K 的 FP32 \texttt{acc}。
\texttt{convert\_acc\_rowcol} 只做布局代数，把 MMA accumulator 重解释为
逻辑 $(\text{row},\text{col})$ 视图；数值仍留在原 fragment 中。

归约写入 \texttt{kReductionBuffers} 个 shared buffer。buffer 少于 warp 时，
各组通过 CTA barrier 串行复用；buffer 与 warp 等数时，每个 warp 直接写一份。
rank32 使用 \texttt{col\^{}row} 与 4 列附加间距，其他 rank 使用一列保护间距。
最后的协作 epilogue 才跨 buffer 求和、折入 R16 scale 并转 BF16：

\begin{lstinputlisting}[linerange={730-745},firstnumber=auto,
  title={quantize\_impl.cuh L730-745（fused epilogue：归约、scale 与单次 BF16 舍入）}]{../../../../cache-dit/csrc/kernels/svdq/cute/fp4/quantize/quantize_impl.cuh}
\end{lstinputlisting}

因此 fused 仍然满足四条简化：没有全局 FP32 LoRA scratch、没有跨 K CTA
atomic、没有 finalize kernel、没有额外激活 repack。$D$ 只在最终写出时
舍入一次，符合低秩分支保持高精度累加的设计目标。

\subsection{split-K：FP32 scratch、relaxed atomics 与 finalize}

split-K 改变的是 $D$ 的归约位置。一个 CTA 覆盖 \texttt{kWarps*32} 行和一个
K128/K256 chunk；当前 \texttt{kWarps=8}，所以行块是 256 行。grid.x 为
$M/256$，grid.y 为 $\lceil K/(64\cdot\texttt{kMaxSlices})\rceil$。每个 warp
仍计算 32 行，但 rank 按 R16 流式处理，只保留一个 R16 的 FP32 accumulator。

部分积按 fragment-major 顺序写入 scratch：
$(\text{row block},\text{rank tile},\text{warp},\text{atom})$ 之后是 32 个
lane，每个 lane 再用 \texttt{red.relaxed.gpu.global.add.v2.f32} 累加相邻两列。
launch 的三个阶段在同一 stream 上依次排列：

\begin{lstinputlisting}[linerange={821-845},firstnumber=auto,
  title={quantize\_impl.cuh L821-845（split-K 的清零、producer 与 finalize）}]{../../../../cache-dit/csrc/kernels/svdq/cute/fp4/quantize/quantize_impl.cuh}
\end{lstinputlisting}

先清零 $M\times R$ 个 FP32 元素，再启动 producer，最后由
\texttt{finalize\_lora\_bf16\_kernel} 解码 scratch。finalize 每线程处理 8 个
相邻 rank 列，用两次 \texttt{float4} 读和一次 \texttt{uint4} 写；R16 scale
在此时折入，每个输出元素只做一次 BF16 舍入。K64 尾部由 predicate 处理，
不要求 $K$ 是 128 或 256 的倍数，只要求公共约定中的 $K\equiv0\pmod{64}$。

数值边界也随之改变：qact/ascales 的每个位置仍只有一个写者，可以保持确定性；
不同 K CTA 对 $D$ 的 relaxed FP32 加法次序不固定，因此 split-K 的 $D$
不保证逐次运行逐位相同，比较时应使用数值误差门槛。于是“无 scratch、
atomic、finalize”只能用于 fused；split-K 有意引入前三者，但仍不需要
额外的激活 repack。

\section{图示}

\begin{figure}[H]
\centering
\begin{tikzpicture}[font=\scriptsize,
  box/.style={draw, rounded corners=2pt, align=center, inner sep=3pt},
  arr/.style={->, thick}]
  % kNative
  \node[anchor=west, font=\small\bfseries] at (0,3.1) {(a) \texttt{kNative}：FP32 fragment pack};
  \node[box, fill=tkOrange!22] (nprod) at (1.6,1.7) {native producer\\（另一个代码族）};
  \node[box, fill=tkFillBlue] (nd) at (5.0,1.7) {$D$：FP32\\MMA-fragment pack};
  \node[box, fill=tkFillBlue] (ng) at (8.6,1.7) {GEMM gather\\gmem$\to$register};
  \node[box, fill=tkPurple!18] (nsc) at (8.6,0.3) {在线乘 \texttt{lora\_scales[i]}};
  \node[box, fill=tkFillGreen] (nep) at (11.9,1.7) {LoRA-up MMA\\必须 mainloop 之后};
  \draw[arr] (nprod) -- (nd);
  \draw[arr] (nd) -- (ng);
  \draw[arr] (ng) -- (nep);
  \draw[arr] (nsc) -- (nep);
  % kRowMajor
  \node[anchor=west, font=\small\bfseries] at (0,-0.9) {(b) \texttt{kRowMajor}：row-major BF16（本章 producer）};
  \node[box, fill=tkGreen!25, text width=2.7cm] (rprod) at (1.7,-2.3) {CuTe producer\\uni / r32 / r64 / r128};
  \node[box, fill=tkGreen!25, text width=2.4cm] (rd) at (5.2,-2.3) {$D$：[M,R] BF16\\scale 已折入};
  \node[box, fill=tkFillGreen] (rtma) at (8.6,-2.2) {TMA 加载\\复用 A/B stage};
  \node[box, fill=tkFillGreen] (rep) at (11.9,-2.2) {LoRA-up MMA\\可 mainloop 之前};
  \draw[arr] (rprod) -- (rd);
  \draw[arr] (rd) -- (rtma);
  \draw[arr] (rtma) -- (rep);
  % 注释
  \node[anchor=north west, align=left, text width=16.2cm, font=\tiny] at (0,-3.1)
    {(a) 的 $D$ 不进 TMA stage，gather 直读 gmem，主循环结束后才能起 LoRA 项；
     (b) 的 $D$ 是普通行主序张量，与权重 tile 一样走 TMA 流水线，stage 复用
     使 LoRA 项可与 NVFP4 主积重叠（第~\ref{ch:40}~章）。两者数学等价
     （\texttt{gemm/gemm\_w4a4.h} L32），差别只在 $D$ 到达 MMA 的路径。};
\end{tikzpicture}
\caption{两种 LoRA 输入 layout 及其下游消费方式。(a) native：FP32 fragment
pack + gmem gather + 在线 scale；(b) row-major：BF16 普通张量 + TMA +
stage 复用。}
\label{fig:39-lora-layout}
\end{figure}

图~\ref{fig:39-lora-layout} 说明，选择 layout 时实际在选择
\textbf{$D$ 的精度、搬运路径与 LoRA 项的调度位置}。\texttt{kNative} 的
per-group scale 使用定长 \texttt{float[8]} 表，因此限 $R\le128$；
\texttt{kRowMajor} 不受这张表限制。不过当前 producer 自身仍只接受到 128
的特定 rank，默认 BN=256/BK=128 GEMM tile 还要求 $R\equiv0\pmod{32}$。
“没有 \texttt{float[8]} 限制”不能扩写成“任意 rank 都可运行”。

\section{性能分析}

当前记录来自 cache-dit 在 2026-11-11 的 RTX PRO 5000 单会话测量。
复测入口为 \texttt{bench/kernels/}\allowbreak\texttt{bench\_svdquant.py}；
5 个 shape 与 rank32/64/128 组成 15 格，warmup 10、iters 50、median，
同机三轮再取 median-of-3。\textbf{本书未复测}，这些数字也不能替代
跨会话稳定验收。

\begin{center}\small
\begin{tabular}{@{}lll@{}}
\toprule
观察层 & 当前记录 & 正确读法 \\
\midrule
producer 算子级 Event & geomean $1.1088\times$ & native/CuTe 延迟比，15/15 快 \\
15 格范围 & $1.037\times$--$1.213\times$ & 不是每种 shape 都获得相同收益 \\
NCU：uni split-K（4 格） & kernel 约 $1.17\times$ & 仍是手写实现 \\
NCU：特化 fused（1 格） & kernel 约 $1.30\times$ & rank32、16K，接近 DRAM 带宽墙 \\
NCU：特化 split-K（10 格） & kernel 约 $1.08\times$ & rank64/128，仍有 shared 冲突 \\
split-K e2e 的 producer 段 & native $94.4\,\mu s$ / CuTe $101.0\,\mu s$ & kernel 更快不保证整段更快 \\
\bottomrule
\end{tabular}
\end{center}

这里不能把 $1.1088\times$ 全部归因于 TiledMMA/TiledCopy。15 格中有 4 格仍走
uni split-K，1 格走特化 fused，10 格走特化 split-K；rank 特化主要扭转了
rank64/128 的旧回退。NCU 行用于解释路径差异，Event 行用于报告算子延迟，
两类指标不能拼接成同一条时间线。split-K 的 finalize 中位约 $2\,\mu s$，
但 e2e 还包括 scratch 分配、清零、launch 与 allocator 上下文，所以最后一行
特意保留了“kernel 快、producer 段总时间反而慢”的反例。

作为历史对照，2026-11-05 的旧手写 fused 实现曾在
$M=4096,N=K=3072,R=32$ 上记录 CuTe $32.19\,\mu s$、native
$31.94\,\mu s$（nsys kernel-only），只能称为当时接近持平。它既不代表当前
split-K 调度，也不能与上面的 Event geomean 逐格比较。早期从单 warp 全 K 到
多 warp 切 K 的优化仍说明并行度的重要性，但旧行主序 qact 写回等数据已不再
对应当前 native pack，本文不再用它们解释现行实现。

\section{实践经验}

\begin{kaogu}[F4]{当前 producer 的优化边界（cache-dit 的调优结论）}
\begin{itemize}[itemsep=0.5pt,topsep=1pt]
  \item 不要为了“统一实现”替换 uni：rank32、$K\le3072$ 的四格仍由它占优；
  \item 不要恢复 \texttt{lora\_act\_to\_bf16} 适配层：当前 producer 直接输出
        row-major BF16 D，额外 launch 与流量只会抵消收益；
  \item 不要把 rank32 fused 的 \texttt{col\^{}row} 归约布局推广到 split-K：
        后者根本不经 shared D buffer；
  \item 不要用 isolated kernel 推断 e2e：scratch、清零、finalize 和 allocator
        上下文必须进入同一条 trace。
\end{itemize}
\end{kaogu}

旧源码曾把 CuTe \texttt{TiledMMA/TiledCopy} 重写列为 TODO；当前
\texttt{quantize\_impl.cuh} 已经落地，而 \textbf{native A/SFA pack 与
row-major BF16 D 的约定没有改变}。稳定的外部约定让实现可以替换，也让
qact bitwise 校验成为回归门。

当前尚未解决的主矛盾转移到特化 split-K：2026-11-11 的 NCU 记录给出约
380.9K 次 shared bank conflict、16.6\% occupancy、198 registers 与
91 KiB smem；native 对照的 conflict 约 1.9K。与此同时，特化 fused 的
DRAM 利用已到 90.2\%。因此后续优化应先定位 split-K 的具体冲突指令，
而不是重新扫描已定下的 rank/M 阈值；这些是 cache-dit 的单会话记录，
本书未复测。

\section{常见坑}

\begin{itemize}[itemsep=1pt,topsep=2pt]
  \item \textbf{S 重复乘除}：量化路径除 \texttt{smooth}，LoRA 路径用原始
        A fragment，因为 \texttt{lora\_down} 已含 $S^{-1}$
        （式~\ref{eq:39-d}）。改动前先确认离线 pack 是否又除了一次 $S$。
  \item \textbf{B fragment 的 even/odd 置换}：通用路径的
        \texttt{load\_lora\_pair} 明示了“相邻 K 不是相邻内存”；特化路径用
        SW32 shared layout 与 \texttt{LDSM\_T} 表达同一转换，不能混用两者的
        索引推导。
  \item \textbf{native pack 不是行主序}：\texttt{local/base} 按 256 行块交织；
        直接按 $[M,K/2]$ 索引 \texttt{qact} 会读错。
  \item \textbf{scale 粒度不是 per-K32}：尽管部分注释这样写，存储大小和 GEMM
        消费约定都是每 16 个 K 元素一个 SFA。
  \item \textbf{SFA 写不能简单合并成 16B}：当前映射每行每个 K64 只写一个
        u32 scale 字；split-K 虽让 CTA 覆盖 256 行，也不会把交织地址变成
        普通连续行。强转 \texttt{uint4} 会写到其他逻辑位置。
  \item \textbf{scratch 就是调度开关}：给 rank 特化入口传有效 FP32 scratch
        必然选择 split-K；要 fused 必须传无效 tensor。只改 grid 或 rank 参数
        不能绕过这条约定。
  \item \textbf{split-K 的 D 不做 bitwise 断言}：跨 K CTA 的 relaxed FP32
        加法次序不固定；qact/ascales 的确定性与 D 的误差门槛要分开验收。
  \item \textbf{shared 与架构边界}：特化 launch 静态限制
        \texttt{kSmemBytes<=101376}，超过默认 48 KiB 时还要 opt-in；producer
        的 host 校验不检查设备，非 SM120 的 device 兜底是 trap。
\end{itemize}

\section{面试要点速查}

\begin{itemize}[itemsep=1pt,topsep=2pt]
  \item Q：为什么量化 producer 里会出现 LoRA-down 的 GEMM？\\
        A：两者共享输入 $X$，同一批原始 A fragment 同时服务
        $\mathcal Q(\widehat X)$ 与 $D$；分开做会增加输入读取和中间张量往返。
  \item Q：LoRA 路径为什么能用原始 $X$ 而不是 $\widehat X$？\\
        A：离线打包时 $\widetilde L_2=L_2S^{-1}$，数学上
        $X\widetilde L_2^{\top}=\widehat XL_2^{\top}$（式~\ref{eq:39-d}）。
  \item Q：为什么有了 rank 特化实现还保留 uni？\\
        A：uni 覆盖 rank 0 与非特化 rank，并在 rank32、$K\le3072$ 的当前
        测量中仍更快；“存在特化实现”不等于“所有 shape 都应选它”。
  \item Q：TiledMMA/TiledCopy 改版改变了哪些约定？\\
        A：改变 A/B 的 shared layout、G2S/S2R copy 与 LoRA fragment 生成方式；
        native qact/ascales pack、原始 A 供 LoRA、row-major BF16 D 都没变。
  \item Q：fused 与 split-K 的核心差别是什么？\\
        A：fused 在 CTA 内用 shared 归约；split-K 跨 K CTA 写 FP32 scratch，
        再 finalize。后者增加中间态与原子累加，也失去 D 的逐位可复现性。
  \item Q：native pack 为什么长成那个样子？\\
        A：它按 SM120 MMA 的 A/SFA fragment 位置交织，让 GEMM 可以 G2R
        直读而不过 TMA/smem；布局由消费者决定。
  \item Q：两种 lora layout 怎么选？\\
        A：row-major BF16 D 可走 TMA 并允许 LoRA-first，是当前 CuTe 默认形态；
        native FP32 fragment D 适合复用既有 native 输出，但只能 main-first。
        具体性能仍按 GPU、shape 与 rank 实测，不能沿用旧 rank128 禁令。
\end{itemize}

\section{最小测试与动手实验}

\begin{lstlisting}[style=console]
# 读者侧验证（cache-dit 仓库，SM120 卡；本书未复跑）：
$ python -m pytest tests/kernels/test_svdquant_cute.py -q
# _require_sm120() 门控：非 SM120 自动 skip（输出 SKIP）
# 覆盖：rank 32/64/128 的 producer 数值、native qact 逐字节一致
#       （torch.equal）、ascales 校验、uni/特化与 fused/split-K 分派、
#       K64 尾块、量化 Linear 的 native/cute parity、非法 M 被拒绝

# CPU 侧观察实验（native pack 映射的复算，numpy）：
# 用 numpy 模拟 base/local 与 writer-lane 偏移，确认按 256 行块交织的
# qact 展开后与 [M,K/2] 行主序重排的 E2M1 码逐元素一致（不需要 GPU）。
# 另用同一组输入比较 fused 与 split-K：qact/ascales 可逐位比，D 用误差门槛。
\end{lstlisting}

\section*{延伸阅读}

\begin{itemize}[itemsep=1pt,topsep=2pt]
  \item nunchaku 的对应 producer：同一融合形态在 SVDQuant 官方引擎里的写法；
  \item SVDQuant 论文（arXiv:2411.05007）Fig.~6：融合原则与
        未融合低秩分支的 57\% 开销测量；
  \item cache-dit \texttt{docs/user\_guide/QUANTIZATION.md}：独立 CuTe 工作流、
        当前 producer dispatch 与读者侧 benchmark 入口；
  \item 本部分另两章：第~\ref{ch:38}~章（数学主线）、第~\ref{ch:40}~章
        （消费本章三种输出的 GEMM）。
\end{itemize}
